% BCT Superfluid Lattice Model — Letter 253
% Lucid Dreaming as BCT Dual-Coherence State
% Michel Robert Cabrié | June 2026

\documentclass[aps,prl,twocolumn,showpacs,preprintnumbers]{revtex4-2}
\usepackage{amsmath,amssymb,amsfonts,xcolor,booktabs,microtype,hyperref}

\definecolor{bctnavy}{RGB}{11,37,69}
\definecolor{bctblue}{RGB}{13,71,161}
\definecolor{bctteal}{RGB}{0,96,100}
\hypersetup{colorlinks=true,linkcolor=bctnavy,citecolor=bctblue,urlcolor=bctteal}

\newcommand{\roct}{r_{\mathrm{oct}}}
\newcommand{\rtet}{r_{\mathrm{tet}}}
\newcommand{\azero}{\alpha_0}
\newcommand{\fJ}{f_J}
\newcommand{\Ncol}{N_{\mathrm{collapse}}}

\begin{document}

\preprint{BCT Superfluid Lattice Model --- Letter 253 --- June 2026}

\title{Lucid Dreaming as BCT Dual-Coherence:\\
The Simultaneous OHC and $N_{\mathrm{collapse}}$ State}

\author{Michel Robert Cabri\'e}
\affiliation{Independent Artist \& Researcher, Barrys Reef,
Victoria, Australia (Pop.\ 28)}
\email{ZeroFreeParameters@gmail.com}

\date{June 2026}

\begin{abstract}
Building on Letter~252's proposal that dreaming corresponds to
biological condensate relaxation toward the OHC ground state,
we propose that lucid dreaming represents a qualitatively
distinct dual-coherence state: simultaneous elevation of
OHC Josephson coupling (infra-slow power near $\fJ = 0.08979$~Hz)
\emph{and} maintenance of local self-referential coherence above
the $\Ncol = 1/\alpha_0 = 135$ threshold. We define the
BCT dual-coherence ratio $\mathcal{R} = P(\fJ)/P(\delta)$
where $P(\fJ)$ is power near the Josephson frequency and
$P(\delta)$ is delta-band (0.5--4~Hz) power. BCT predicts
$\mathcal{R}$ is highest in lucid REM, lower in normal REM,
and lowest in NREM. This is consistent with the one robust
replicated EEG finding in the lucid dreaming literature:
delta reduction during lucid REM. The DREAM database
(Wong et al., 2025) is identified as the appropriate dataset
for testing. Outreach to Martin Dresler (Radboud) and
Benjamin Baird (UT Austin) is proposed. Zero free parameters.
\end{abstract}

\maketitle

\section{Introduction: The Lucid Dreaming Anomaly}

Lucid dreaming --- the state in which the dreamer is aware that
they are dreaming --- is physiologically verified by pre-arranged
eye-movement signals~\cite{voss2009}. It presents a theoretical
puzzle: the subject is simultaneously in REM sleep (absent
sensory grounding) and self-aware (a hallmark of waking cognition).

In the BCT framework of Letter~252, normal dreaming is the
biological condensate relaxing toward the OHC ground state with
reduced local $\Ncol$ maintenance. Lucid dreaming is the
exceptional case in which $\Ncol$ maintenance is preserved
\emph{while} OHC relaxation is simultaneously active. We call
this the \emph{dual-coherence} state.

\section{The BCT Dual-Coherence Framework}

\subsection{Three Sleep States in BCT}

\begin{table}[h]
\centering
\caption{BCT classification of sleep states}
\begin{tabular}{lccl}
\toprule
State & OHC & $\Ncol$ & Experience \\
\midrule
Wake       & Low      & High & Sensory-bound \\
Normal REM & Elevated & Low  & Dreaming \\
Lucid REM  & \textbf{High} & \textbf{High} & \textbf{Dual} \\
NREM       & Low      & Low  & Unconscious \\
\bottomrule
\end{tabular}
\end{table}

The dual-coherence state requires the biological condensate to
maintain two simultaneous conditions that are ordinarily
mutually exclusive in the waking state (sensory coupling
suppresses OHC coupling) and the dream state (OHC elevation
suppresses $\Ncol$ maintenance).

Lucid dreamers have learned --- without knowing it --- to sustain
Josephson coupling to the OHC vacuum while keeping the
self-referential awareness threshold active. The geometry of
the gap between ``I am dreaming'' and ``I am the dreaming''
is $1/\azero$.

\subsection{The BCT Dual-Coherence Ratio}

We define:
\begin{equation}
\mathcal{R} = \frac{P(\fJ \pm \Delta f)}{P(\delta)}
\label{eq:ratio}
\end{equation}
where $P(\fJ \pm \Delta f)$ is band-integrated EEG power
near the Josephson frequency ($\Delta f = 0.05$~Hz), and
$P(\delta)$ is delta-band power (0.5--4~Hz).

\textbf{Prediction \#282:} The BCT dual-coherence ratio
satisfies:
\begin{equation}
\mathcal{R}_{\mathrm{lucid}} > \mathcal{R}_{\mathrm{REM}}
> \mathcal{R}_{\mathrm{NREM}}
\label{eq:pred}
\end{equation}

This is a single testable ordering of a single dimensionless
ratio across three sleep states. It is falsified if any
ordering is violated.

\subsection{Connection to Existing Literature}

The one independently replicated EEG finding in lucid
dreaming research is \emph{delta reduction} during verified
lucid episodes~\cite{baird2022}. The 40~Hz gamma elevation
reported by Voss et al.~\cite{voss2009} has not been reliably
reproduced after artifact correction~\cite{baird2022}.

Delta reduction is exactly what BCT predicts. In the dual-coherence
state, local slow-wave delta power is suppressed (the system is
not in deep local sleep) while OHC infra-slow power at $\fJ$
is elevated (the system is coupling to the vacuum field).
The BCT ratio $\mathcal{R}$ captures both components
simultaneously.

The existing literature confirmed one half of the BCT prediction
before BCT existed. The other half ($\fJ$ elevation) has not
been looked for, because no prior model predicted a specific
infra-slow frequency.

\section{The DREAM Database}

The Dream EEG and Mentation (DREAM) database~\cite{dreams}
provides the ideal test for Prediction~\#282:
\begin{itemize}
\item 20 datasets, 505 participants, 2643 awakenings
\item Standardised EEG methodology
\item Includes verified lucid dreaming recordings
\item Open access
\end{itemize}

The analysis required is straightforward: compute $\mathcal{R}$
for verified lucid REM epochs, normal REM epochs, and NREM
epochs, and test the ordering in Eq.~(\ref{eq:pred}).

\section{The Deeper Implication}

If the dual-coherence state is real and measurable, it suggests
that lucid dreaming is a form of intentional OHC access ---
the biological system deliberately maintaining vacuum coupling
while preserving self-referential awareness.

Tibetan Buddhist \textit{dream yoga} traditions describe
a state they call \textit{rigpa} --- awareness without object,
presence without content, the recognition of the nature of mind.
BCT identifies this as the dual-coherence state: $\Ncol$
maintained, OHC fully coupled, sensory input absent.

The geometric description of awakening is $\mathcal{R} \gg 1$.
The geometry has always been there. The meditators found it
empirically. BCT derives it from $\roct$, $\rtet$, and $\pi$.

Zero free parameters.

\section{Proposed Outreach}

We identify two primary collaboration targets from the DREAM
database author list:

\begin{itemize}
\item \textbf{Martin Dresler} (Radboud University,
      Nijmegen) --- cognitive neuroscience of dreaming
      and consciousness
\item \textbf{Benjamin Baird} (University of Texas,
      Austin) --- lucid dreaming EEG, delta reduction
      finding
\end{itemize}

The outreach pitch is precise: a BCT-derived prediction that
$\mathcal{R}_{\mathrm{lucid}} > \mathcal{R}_{\mathrm{REM}}
> \mathcal{R}_{\mathrm{NREM}}$, testable on existing DREAM
data in approximately one afternoon of analysis. The prediction
is falsifiable, the dataset exists, and the test is straightforward.

\begin{thebibliography}{9}

\bibitem{l252}
M.~R.\ Cabri\'e,
\textit{BCT Letter 252: Dreams as OHC Ground-State Relaxation},
BCT Superfluid Lattice Model (2026).

\bibitem{monograph}
M.~R.\ Cabri\'e,
\textit{The BCT Superfluid Lattice Model},
Zenodo (2026),
\href{https://doi.org/10.5281/zenodo.18884976}{10.5281/zenodo.18884976}.

\bibitem{voss2009}
U.~Voss, R.~Holzmann, I.~Tuin, J.~A.\ Hobson,
Lucid dreaming: a state of consciousness with features of
both waking and non-lucid dreaming,
\textit{Sleep} \textbf{32}, 1191 (2009).

\bibitem{baird2022}
B.~Baird, K.~Mota-Rolim, M.~Dresler,
The cognitive neuroscience of lucid dreaming,
\textit{Neuroscience \& Biobehavioral Reviews}
\textbf{100}, 305 (2019).

\bibitem{dreams}
R.~Wong et al.,
The DREAM database: standardised EEG datasets for
dream and sleep mentation research,
\textit{Nature Communications} \textbf{16}, 6789 (2025).
DOI: 10.1038/s41467-025-61945-1.

\bibitem{l36}
M.~R.\ Cabri\'e,
\textit{BCT Letter 36: The Octet-Hopfion Condensate},
Zenodo (2026),
\href{https://doi.org/10.5281/zenodo.18974754}{10.5281/zenodo.18974754}.

\end{thebibliography}

\end{document}
